% Figure 3. Physical Hilbert space. What the reader must see: relabelling the argument matches permuting the spin
% factors, up to the statistics sign; parity is a separate character; inversion flips handedness, which no rotation
% does. Drawing: X, sigma X and E*X in a row, the spin slots in label order with leaders from the digits, so the
% factors visibly follow the labels; the inversion triad on a sphere after Harter (1993) fig. 5.1.1.
% Colours: oxygen red only; gold unused. Hatch unused. Digits are the column labels.
\documentclass[10pt,border=1mm]{standalone}
\usepackage[T1]{fontenc}
\usepackage{figurestyle}
\input{primitives.tex}
\input{molecules.tex}
\input{numbers.tex}
\begin{document}
\begin{tikzpicture}[plate,slot/.style={draw=PlateInk,line width=.45pt,minimum width=.62cm,minimum height=.48cm,inner sep=1pt,fill=white,font=\fontsize{8}{9}\selectfont}]
\path[use as bounding box] (0,0) rectangle (15,9.0);
% ------------------------------------------------ X, sigma X, E*X with the spin slots
\foreach \x/\la/\lb/\pname/\sa/\sb/\cap in {2.0/1/2/mol3d-water-X/1/2/{$X$},6.4/2/1/mol3d-water-X/2/1/{$\sigma X=XP_\sigma$},10.8/1/2/mol3d-water-minusX/1/2/{$E^*X=-X$}} {
  \begin{scope}[shift={(\x,7.1)}]\molsolid\def\molA{\la}\def\molB{\lb}\pic[scale=1.0] at (0,0) {\pname};\end{scope}
  \node[lab] at (\x,8.45) {\cap};
  \node[slot] (s1) at (\x-.9,4.7) {$\xi_{\sa}$}; \node at (\x-.45,4.7) {$\otimes$};
  \node[slot] (s2) at (\x,4.7) {$\xi_{\sb}$}; \node at (\x+.45,4.7) {$\otimes$};
  \node[slot] (s3) at (\x+.9,4.7) {$\xi_3$};
  \node[sub,anchor=north] at (\x-.9,4.4) {slot $1$}; \node[sub,anchor=north] at (\x,4.4) {slot $2$}; \node[sub,anchor=north] at (\x+.9,4.4) {slot $3$};}
% leaders: slot i to the nucleus carrying digit i; atom positions from the pics (H_A, H_B, O at pic scale 1), molecule centre at y = 7.1
% X: digit 1 on H_A (left), digit 2 on H_B (right).  sigma X: digit 2 on H_A, digit 1 on H_B.  E*X: positions negated, digit 1 on -H_A (right, low).
\draw[leader] (2.0-.9,4.95)--(2.0-.93,6.3);   \draw[leader] (2.0,4.95)--(2.0+.65,6.5);   \draw[leader] (2.0+.9,4.95)--(2.0+.3,6.75);
\draw[leader] (6.4-.9,4.95)--(6.4+.65,6.5);   \draw[leader] (6.4,4.95)--(6.4-.93,6.3);   \draw[leader] (6.4+.9,4.95)--(6.4+.3,6.75);
\draw[leader] (10.8-.9,4.95)--(10.8+.93,6.3); \draw[leader] (10.8,4.95)--(10.8-.65,6.5); \draw[leader] (10.8+.9,4.95)--(10.8+.2,7.0);
\draw[harrow] (3.5,7.1)--(4.9,7.1); \node[lab,anchor=south] at (4.2,7.2) {$\sigma=(12)$}; \node[sub,anchor=north] at (4.2,7.0) {relabel};
\draw[harrow] (7.9,7.1)--(9.3,7.1); \node[lab,anchor=south] at (8.6,7.2) {$E^*$}; \node[sub,anchor=north] at (8.6,7.0) {invert};
\node[sub,anchor=north,text width=4.2cm,align=center] at (6.4,4.0) {the same points of space; slot $i$ belongs to the nucleus carrying digit $i$, so the factors follow the labels};
\node[sub,anchor=north,text width=4.2cm,align=center] at (10.8,4.0) {planar, so $-X$ is also a rotated copy; the triad shows what inversion does};
\node[sub,anchor=north,text width=4.2cm,align=center] at (2.0,4.0) {a product value $\Psi(X)=\xi_1\otimes\xi_2\otimes\xi_3$; sums of such values by linearity};
% ------------------------------------------------ the two relations and the statistics sign
\node[lab,anchor=west] at (.4,2.1) {$\Psi(\sigma X)\sim\stat(\sigma)\,\sigma\act\Psi(X)$,\quad $\stat(\sigma)=\prod_\alpha(\operatorname{sgn}\sigma_\alpha)^{2I_\alpha}$:\ $-1$ for ${}^1\mathrm H$, $+1$ for ${}^2\mathrm H$, ${}^{16}\mathrm O$ not permuted};
\node[lab,anchor=west] at (.4,1.4) {$\Psi(E^*X)\sim\pm\Psi(X)$ on $\hilb^\pm$:\quad parity is a separate character, $\chi_\pm(\sigma^*)=\pm1$, and $\stat$ ignores the star};
\node[sub,anchor=west] at (.4,.75) {equalities hold almost everywhere on $\conf$; the physical space $\hilb\subseteq\hilb_{\mathrm{ambient}}$ is cut out by these two conditions};
% ------------------------------------------------ inversion triad on a sphere
\begin{scope}[shift={(13.6,7.1)}]
  \hatom{0}{0}{.72}{PlateInk!6}
  \draw[harrow,line width=.6pt] (0,0)--(0,1.0); \node[sub,anchor=south] at (0,1.02) {$z$};
  \draw[harrow,line width=.6pt] (0,0)--(1.0,0); \node[sub,anchor=west] at (1.02,0) {$y$};
  \draw[harrow,line width=.6pt] (0,0)--(-.48,-.58); \node[sub,anchor=north east,inner sep=1pt] at (-.48,-.58) {$x$};
  \draw[hidden,-{Stealth[length=1.6mm]}] (0,0)--(0,-.95); \draw[hidden,-{Stealth[length=1.6mm]}] (0,0)--(-.95,0); \draw[hidden,-{Stealth[length=1.6mm]}] (0,0)--(.46,.56);
  \node[sub,anchor=north,text width=2.6cm,align=center] at (0,-1.15) {$E^*$ sends the triad to its negative (dashed): right-handed becomes left-handed};
\end{scope}
\end{tikzpicture}
\end{document}
